% Bitwise Hacks

High-performance bitwise manipulations. $O(1)$ submask traversal and hardware-accelerated queries.

\begin{lstlisting}
// 1. Core Bit Operations
x | (1LL << i) // Set i-th bit to 1
x & ~(1LL << i) // Set i-th bit to 0
x ^ (1LL << i) // Toggle i-th bit
(x >> i) & 1LL // Check if i-th bit is 1
long long lsb = x & -x; // Isolate lowest set bit (Fenwick)
long long clear_lsb = x & (x - 1); // Clear lowest set bit

// 2. Bulk & Range Bit Manipulations (k bits, Range [l, r])
(1LL << k) - 1 // Mask of k lowest bits (000...111)
x ^ ((1LL << k) - 1) // Toggle k lowest bits of x
((1LL << (r - l + 1)) - 1) << l // Mask for range [l, r] (inclusive)
x ^ (((1LL << (r - l + 1)) - 1) << l) // Toggle all bits in range [l, r]
x | (x + 1) // Set the lowest unset (0) bit to 1

// 3. Submask Traversal
for (int sub = mask; sub > 0; sub = (sub - 1) & mask) {} // Iterate all submasks

// 4. Built-ins & Helper Functions
__builtin_popcountll(x) // Count set bits
__builtin_ctzll(x) // Count trailing zeros (x > 0)
__builtin_clzll(x) // Count leading zeros (x > 0)
__lg(x) // Fast Floor(log2(x)) -> Sparse Table query
long long next_p2(long long x) { return x <= 1 ? 1 : 1LL << (64 - __builtin_clzll(x - 1)); } // Next power of 2
long long safe_mod_mul(unsigned long long a, unsigned long long b, unsigned long long m) { // O(1) mul modulo without overflow
    long long res; return __builtin_mul_overflow(a, b, &res) ? (long long)(((unsigned __int128)a * b) % m) : res % m;
}
\end{lstlisting}
